\section{Conclusion: recurrence moves into the trace fiber}
\label{sec:conclusion}

SD-C10 restores recurrence to the tensor-prime base grammar.  Every adjacent
pair supports a mixed two-cycle, yet the universal positive cocycle prevents
that cycle from closing in the identity fiber.  The canonical trace then
recovers exactly the old atom-loop ledger and its Euler determinant.  In this
precise sense, Paper07's acyclicity has moved from the base graph into the
cocycle's positive cone.

The gain is real but bounded.  The base periodic structure is richer, the
transfer is noncommutative, and its chiral fourth trace moves with height.
Nevertheless self-adjointness forces inverse closure.  Every edge then meets
its adjoint, $gg^{-1}$ becomes visible, and the resulting positive quadratic
trace diverges.  The first determinant that exists is $\det_3$, which has
already subtracted the very recurrent term one hoped to retain.

The stage therefore gives a sharp trilemma:
\[
 \boxed{
 \begin{gathered}
 \text{positive-cone analytic trace: exact Euler ledger},\\
 \text{inverse-closed chiral trace: positive recurrent backtracking},\\
 \text{available regularization: deletes the first backtracking term}.
 \end{gathered}}
\]
Fuglede--Kadison and Brown data do not close the triangle because they trade
holomorphic divisor information for magnitude geometry.

Paper09 should not search for another fitted phase.  It should test a
holomorphic doubled free-semigroup transfer using separate positive
alphabets for $s$ and $1-s$.  The target is one relative analytic determinant
whose first mixed coefficient survives in a common Schatten domain.  Failure
should be upgraded to an inverse-closure/ideal obstruction theorem.

No external geometric system is developed.  If local inverse phases need a
geometric origin and an additional sign-reversing law, that possibility is a
\texttt{ROUND2\_CLUE}, not part of SD-C10.
